\chapter{Данные: тома и привязки}

Записываемый слой контейнера исчезает вместе с контейнером. Это
правильное поведение: контейнеры должны быть взаимозаменяемыми,
а состояние~--- жить отдельно. В главе разобраны три способа подключить
данные~--- тома, привязки каталогов и \cmd{tmpfs},~--- а также резервное
копирование, права доступа и типичные ошибки при работе с базами данных.

\section{Почему данные нельзя хранить в контейнере}

\begin{code}[language=bash]
$ docker run --rm -it alpine sh
/ # echo "важные данные" > /data.txt
/ # exit
$ docker run --rm -it alpine cat /data.txt
cat: can't open '/data.txt': No such file or directory
\end{code}

Второй контейнер создан из того же образа, но его записываемый слой
пуст. Данные первого контейнера исчезли вместе с ним. Даже если
контейнер не удалять, обновление версии приложения обычно означает
создание нового контейнера из нового образа~--- и снова чистое
состояние.

\begin{defbox}[title={Определение: том}]
\textbf{Том} (\emph{volume})~--- хранилище, управляемое Docker, которое
существует независимо от контейнеров и монтируется в их файловую
систему. Тома живут в каталоге данных Docker (обычно
\file{/var/lib/docker/volumes}) и удаляются только явно.
\end{defbox}

\section{Три способа подключения данных}

\begin{figure}[htbp]
\centering
\begin{tikzpicture}[
  font=\sffamily\scriptsize,
  host/.style={draw=accent!55!black, fill=accent!8, rounded corners=2pt,
    minimum width=3.6cm, minimum height=1.0cm, align=center},
  cont/.style={draw=good!60!black, fill=good!8, rounded corners=2pt,
    minimum width=3.4cm, minimum height=1.0cm, align=center},
  arr/.style={-Latex, thick, accent!70},
  lbl/.style={font=\sffamily\scriptsize, text=muted},
]
  % bind
  \node[host] (h1) at (0,2.2) {каталог хоста\\\cmd{/home/user/site}};
  \node[cont] (c1) at (7.0,2.2) {контейнер\\\cmd{/usr/share/nginx/html}};
  \draw[arr] (h1) -- node[above, lbl] {bind mount} node[below, lbl] {вы сами управляете} (c1);

  % volume
  \node[host] (h2) at (0,0.7) {том \cmd{pgdata}, управляет Docker\\\cmd{/var/lib/docker/volumes/...}};
  \node[cont] (c2) at (7.0,0.7) {контейнер\\\cmd{/var/lib/postgresql/data}};
  \draw[arr] (h2) -- node[above, lbl] {volume} node[below, lbl] {данные переживают контейнер} (c2);

  % tmpfs
  \node[host] (h3) at (0,-0.8) {оперативная память\\ядра};
  \node[cont] (c3) at (7.0,-0.8) {контейнер\\\cmd{/tmp}};
  \draw[arr] (h3) -- node[above, lbl] {tmpfs} node[below, lbl] {исчезает при остановке} (c3);
\end{tikzpicture}
\caption{Три способа подключения данных: привязка каталога, том Docker
и tmpfs в памяти}\label{fig:mounts}
\end{figure}

\begin{center}
\small
\begin{tabularx}{\textwidth}{@{}lXXX@{}}
\toprule
 & \textbf{Том} & \textbf{Привязка} & \textbf{tmpfs} \\
\midrule
Где хранится & каталог данных Docker & любой путь хоста & память \\
Кто управляет & Docker & вы & ядро \\
Живёт после удаления контейнера & да & да & нет \\
Резервное копирование & \cmd{docker run tar} или приложение & обычные средства
хоста & не требуется \\
Типичное применение & базы данных, состояние сервисов & конфигурации,
исходный код при разработке & кэш, временные файлы, секреты \\
Скорость & высокая & зависит от ФС (на macOS медленнее) & очень высокая \\
\bottomrule
\end{tabularx}
\end{center}

\section{Именованные и анонимные тома}

\begin{code}[language=bash]
$ docker volume create pgdata
$ docker volume ls
DRIVER    VOLUME NAME
local     pgdata
$ docker volume inspect pgdata
[
    {
        "Name": "pgdata",
        "Driver": "local",
        "Mountpoint": "/var/lib/docker/volumes/pgdata/_data",
        ...
    }
]
\end{code}

Подключение:

\begin{code}[language=bash]
$ docker run -d --name db \
    -e POSTGRES_PASSWORD=secret \
    -v pgdata:/var/lib/postgresql/data \
    postgres:17

$ docker inspect db --format '{{json .Mounts}}'
[{"Type":"volume","Name":"pgdata","Destination":"/var/lib/postgresql/data",...}]
\end{code}

Если указать только путь внутри контейнера (\flag{-v /var/lib/data}),
Docker создаст \emph{анонимный} том со случайным именем. Он переживёт
перезапуск контейнера, но его легко потерять: \cmd{docker rm -v}
удаляет анонимные тома вместе с контейнером, а \cmd{docker volume prune}
убирает все неиспользуемые.

\begin{code}[language=bash]
$ docker run --rm -v /var/lib/data alpine ls /var/lib/data
$ docker volume ls -f dangling=true
$ docker volume rm pgdata            # только если том не используется
$ docker volume prune                # удалить все неиспользуемые тома
\end{code}

\begin{warnbox}
\cmd{docker volume prune} не спрашивает подтверждения в некоторых
версиях и удаляет данные безвозвратно. Убедитесь, что среди
неиспользуемых томов нет резервных копий или баз, к которым вы
планируете вернуться.
\end{warnbox}

\section{Привязки каталогов}

Привязка (\emph{bind mount}) монтирует каталог или файл хоста внутрь
контейнера. Это основной инструмент разработки: изменения в файлах
сразу видны приложению.

\begin{code}[language=bash]
$ docker run -d --name web -p 8080:80 \
    --mount type=bind,src="$(pwd)/site",dst=/usr/share/nginx/html,readonly \
    nginx:1.27-alpine

# эквивалентная краткая запись
$ docker run -d --name web2 -p 8081:80 \
    -v "$(pwd)/site:/usr/share/nginx/html:ro" \
    nginx:1.27-alpine
\end{code}

\begin{notebox}
\cmd{-v} и \cmd{--mount} решают одну задачу, но \cmd{--mount}
явнее и безопаснее: если источник не существует, он завершится
с ошибкой, а \cmd{-v} молча создаст каталог на хосте от имени root.
В скриптах и документации предпочитайте \cmd{--mount}.
\end{notebox}

Права и SELinux:

\begin{itemize}
  \item \cmd{:ro}~--- только чтение;
  \item \cmd{:z}~--- перемаркировать каталог для SELinux с разделением
        между контейнерами;
  \item \cmd{:Z}~--- эксклюзивная метка для одного контейнера.
\end{itemize}

\begin{warnbox}
Права доступа внутри контейнера определяются UID/GID процесса,
а не именем пользователя. Если приложение работает от UID 10001,
а каталог на хосте принадлежит UID 1000, запись не удастся.
Проверяйте \cmd{ls -ln}, а не \cmd{ls -l}, и выравнивайте UID.
\end{warnbox}

\section{tmpfs}

Временное хранилище в оперативной памяти. Подходит для кэша, сокетов,
временных файлов и данных, которые не должны попасть на диск.

\begin{code}[language=bash]
$ docker run --rm -it --tmpfs /tmp:size=64m alpine sh
$ docker run -d \
    --mount type=tmpfs,dst=/run/app,tmpfs-size=32m,tmpfs-mode=1770 \
    myapp:1.0
\end{code}

Содержимое \cmd{tmpfs} исчезает при остановке контейнера и занимает
память хоста, учитываемую в лимите \flag{--memory}.

\section{Практика: PostgreSQL с томом}

Запустим базу с именованным томом, создадим данные, перезапустим
контейнер и убедимся, что они сохранились.

\begin{code}[language=bash]
$ docker run -d --name db \
    -e POSTGRES_PASSWORD=secret \
    -e POSTGRES_DB=app \
    -p 127.0.0.1:5432:5432 \
    -v pgdata:/var/lib/postgresql/data \
    postgres:17

$ docker exec -it db psql -U postgres -d app -c \
    "CREATE TABLE notes(id serial primary key, body text);"
$ docker exec -it db psql -U postgres -d app -c \
    "INSERT INTO notes(body) VALUES ('тома работают');"

$ docker restart db
$ docker exec -it db psql -U postgres -d app -c \
    "SELECT * FROM notes;"
 id |      body
----+----------------
  1 | тома работают
(1 row)
\end{code}

Логическая резервная копия средствами PostgreSQL:

\begin{code}[language=bash]
$ docker exec -t db pg_dump -U postgres app > backup.sql
$ docker exec -i db psql -U postgres app < backup.sql
\end{code}

Такая копия переносима между версиями и удобна для небольших баз.
Для больших баз используют физические бэкапы и инструменты вроде
pgBackRest, запуская их в отдельном контейнере с тем же томом.

\section{Резервное копирование томов}

Универсальный способ~--- подключить том и каталог хоста к временному
контейнеру и упаковать данные.

\begin{code}[language=bash]
$ docker run --rm \
    -v pgdata:/data \
    -v "$(pwd)":/backup \
    alpine \
    tar czf /backup/pgdata.tar.gz -C /data .

# восстановление
$ docker run --rm \
    -v pgdata:/data \
    -v "$(pwd)":/backup \
    alpine \
    sh -c "cd /data && tar xzf /backup/pgdata.tar.gz"
\end{code}

\begin{warnbox}
Не копируйте каталог данных работающей базы «на горячую» без
согласованного бэкапа: файлы могут быть в несогласованном состоянии.
Останавливайте базу, используйте \cmd{pg\_dump} или штатные средства
резервного копирования.
\end{warnbox}

\section{Тонкости, о которых узнают поздно}

\subsection{VOLUME в Dockerfile}

Инструкция \cmd{VOLUME} объявляет точку монтирования. Данные,
записанные в этот путь \emph{после} инструкции во время сборки,
не сохраняются в образе: при запуске поверх будет пустой том.

\begin{code}[language=dockerfile]
VOLUME /var/lib/data
RUN echo "данные" > /var/lib/data/file.txt   # этот файл не попадёт в образ
\end{code}

Поэтому данные приложения обычно готовят до \cmd{VOLUME} или вообще
не объявляют том в Dockerfile, а задают его при запуске.

\subsection{Права на именованный том}

При первом подключении пустого именованного тома Docker копирует
в него содержимое каталога из образа, включая владельца и права.
Если каталог в образе принадлежит нужному UID, проблема решается
сама. В противном случае используйте init-контейнер или
\cmd{--user} с подходящим UID.

\begin{code}[language=bash]
$ docker run --rm -v pgdata:/data alpine chown -R 10001:10001 /data
\end{code}

\subsection{Копирование файлов без томов}

Иногда достаточно \cmd{docker cp}: скопировать конфигурацию в контейнер
или вытащить логи и артефакты.

\begin{code}[language=bash]
$ docker cp ./nginx.conf web:/etc/nginx/conf.d/default.conf
$ docker cp web:/var/log/nginx/access.log ./access.log
\end{code}

\begin{summarybox}
\begin{itemize}
  \item Записываемый слой эфемерен; состояние выносится в тома
        или привязки.
  \item Том управляется Docker и живёт до явного удаления; привязка~
        это каталог хоста, которым управляете вы.
  \item \cmd{tmpfs} хранит данные в памяти и исчезает при остановке.
  \item Права определяются UID/GID, а не именами пользователей.
  \item Резервное копирование томов делается через временный контейнер
        с \cmd{tar}; базы данных лучше копировать их собственными
        средствами.
\end{itemize}
\end{summarybox}

\section*{Вопросы для самопроверки}
\begin{selfcheck}
  \item Почему данные, записанные в контейнер, теряются при его удалении?
  \item Чем именованный том отличается от анонимного?
  \item Когда выбрать привязку каталога, а когда~--- том?
  \item Что произойдёт с файлами, созданными в каталоге \cmd{VOLUME}
        во время сборки?
  \item Как сделать резервную копию тома и восстановить её?
  \item Почему приложение не может писать в привязанный каталог,
        хотя «права выставлены»?
\end{selfcheck}
